\chapter{Informe de sostenibilidad}

\section{Autoevaluacion}
La sostenibilidad de Odín se analiza desde tres dimensiones: económica, ambiental y social. El proyecto no trata la sostenibilidad como una sección añadida al final, sino como una consecuencia directa de su diseño: sustituir dependencias externas por infraestructura propia cambia los costes, el consumo, la privacidad y la responsabilidad del usuario.

\section{Dimension ambiental}
La dimensión ambiental se centra en el equilibrio entre ejecutar IA local y evitar el uso constante de infraestructuras remotas. El diseño debe minimizar consumo base mediante servicios ligeros, activar la GPU solo cuando la carga lo justifique y reutilizar el hardware para otros fines al final del ciclo principal del proyecto.

\section{Dimension economica}
La viabilidad económica se basa en pasar de un modelo OPEX de suscripciones y APIs externas a un modelo CAPEX de infraestructura propia. La memoria deberá comparar esta inversión con alternativas comerciales, teniendo en cuenta no solo euros, sino independencia, aprendizaje y control.

\section{Dimension social}
Socialmente, Odín responde a una preocupación creciente: la pérdida de control sobre información íntima en entornos digitales opacos. Su aportación es demostrar que una parte de la IA personal puede ejecutarse de forma privada y auditable. Al mismo tiempo, se reconocerá la barrera técnica del autoalojamiento, que limita su adopción por usuarios no especializados.

\section{Matriz de sostenibilidad}
\placeholder{Incluir tabla final con impactos positivos, impactos negativos, métricas observadas y medidas de mejora.}
